% Part: second-order-logic
% Chapter: sol-and-sets
% Section: cardinalities

\documentclass[../../../include/open-logic-section]{subfiles}

\begin{document}

\olfileid[tr]{sol}{set}{crd}

\olsection{Kümelerin Kardinaliteleri}

\begin{explain}
Tanım kümesinin sonlu ya da sonsuz, !!{enumerable} ya da !!{nonenumerable}
olduğunu ifade edebildiğimiz gibi, $\Domain{M}$'nin bir alt kümesinin sonlu
ya da sonsuz, !!{enumerable} ya da !!{nonenumerable} olma özelliğini de
tanımlayabiliriz.
\end{explain}

\begin{prop}
$\fn{Inf}(X)\ident$
\begin{multline*}
\lexists[u][(\lforall[x][\lforall[y][(\eq[u(x)][u(y)] \lif 
      \eq[x][y])]] \land {} \\\\lexists[y][(X(y) \land \lforall[x][(X(x)
      \lif \eq/[y][u(x)])]])]
\end{multline*}
formülü, bir değişken ataması $s$ bakımından ancak ve ancak $s(X)$ sonsuzsa
sağlanır.
\end{prop}

\begin{prop}
$\fn{Count}(X)\ident$
\begin{multline*}
\lexists[z][\lexists[u][(X(z) \land 
    \lforall[x][(X(x) \lif X(u(x)))] \land {} \\ \lforall[Y][((Y(z) \land
      \lforall[x][(Y(x) \lif Y(u(x)))]) \lif X = Y])]])
\end{multline*}
formülü, bir değişken ataması $s$ bakımından ancak ve ancak $s(X)$
!!{enumerable} ise sağlanır.
\end{prop}

Cantor Teoreminden !!{nonenumerable} kümeler bulunduğunu, hatta sonsuz
büyüklüklerin sonsuz sayıda farklı düzeyi olduğunu biliyoruz. Küme teorisi,
küme büyüklüklerinin bütün bir aritmetiğini geliştirir ve kümelere sonsuz
kardinal sayılar atar. Doğal sayılar, sonlu kümelerin büyüklüklerini ölçen
kardinal sayılardır. !!{denumerable} kümelerin kardinalitesi, $\aleph_0$
(``alef-sıfır'') denen ilk sonsuz kardinalitedir. Sonraki sonsuz büyüklük
$\aleph_1$'dir. Bu, bir kümenin sayılabilir (yani $\aleph_0$ büyüklüğünde)
olmadan sahip olabileceği en küçük büyüklüktür. ``$X$, $\aleph_0$
büyüklüğündedir'' ifadesini
$\fn{Aleph}_0(X)\liff\fn{Inf}(X)\land\fn{Count}(X)$ olarak tanımlayabiliriz.
$X$, ancak ve ancak bütün alt kümeleri sonlu ya da $\aleph_0$ büyüklüğünde,
ama kendisi $\aleph_0$ büyüklüğünde değilse $\aleph_1$ büyüklüğündedir.
Dolayısıyla bunu
$\fn{Aleph_1}(X)\ident\lforall[Y][(Y\subseteq X\lif
(\lnot\fn{Inf}(Y)\lor\fn{Aleph}_0(Y)))]\land\lnot\fn{Aleph}_0(X)$
formülüyle ifade edebiliriz. $\aleph_2$ büyüklüğünde olma benzer biçimde
tanımlanır ve bu böyle sürer.

Özel önem taşıyan bir büyüklük, kontinuumun kardinalitesi denen büyüklüktür.
Bu, $\Pow{\Nat}$'ın, eşdeğer olarak da $\Real$'in büyüklüğüdür. Bir kümenin
kontinuum büyüklüğünde olduğu da ikinci dereceden mantıkta ifade edilebilir,
fakat biraz daha çok çalışma gerektirir.

\end{document}
